\chapter{Conclusion and Future Work}\label{chap:conclusion}

\section{Summary of the Work}

This project set out to build a platform that shortens the path from an emergency to a properly equipped response while
remaining safe, private and trustworthy. \ProjectName{} delivers that platform: a Flutter mobile app and a Next.js web console,
backed by a Spring Boot API with a PostgreSQL database, a separate AI model service, and managed services for identity,
storage, messaging, identity verification and maps. It serves ten roles through 97 functional requirements and 64 use cases,
and its design is documented by 41 diagrams, a 26-table data dictionary and a traceability matrix.

A citizen can report an emergency in one action, even offline; the platform classifies the report and flags duplicates; a
dispatcher assigns the nearest suitable unit from a live queue; the crew works the mission from a phone; an ambulance crew
sends the patient's vital signs to the receiving hospital in advance; and the reporter follows every step. Underneath, a
small set of well-chosen patterns keeps this reliable and private.

\section{Achievement of the Objectives}

\Cref{tab:objectives} relates each objective of \cref{chap:introduction} to what was delivered and where the evidence is.

\begin{table}[!htbp]
\caption{Objectives and the evidence that they were met.}\label{tab:objectives}
\centering\tabfont\rowband
\begin{tabular}{@{}>{\raggedright\arraybackslash}p{4.3cm}>{\raggedright\arraybackslash}p{6.4cm}>{\raggedright\arraybackslash}p{3.5cm}@{}}
\rowcolor{najdanavy}\textcolor{white}{\sffamily\bfseries\small Objective} & \textcolor{white}{\sffamily\bfseries\small Delivered} & \textcolor{white}{\sffamily\bfseries\small Evidence} \\
Shorten the time to assignment with a prioritised, de-duplicated live queue & Live queue with AI priority and duplicate suggestion; unit ranking by distance and ETA; urgent notifications & \cref{sec:fr}; UC-29 to UC-32; TC-29 \\
Safe, fast reporting under stress & One-tap SOS, automatic location, idempotent submission, durable offline outbox & UC-20, UC-28; TC-08 to TC-11 \\
Role-appropriate visibility, including audited medical data & Role and ownership checks; consent- and role-gated medical profile with audit & \cref{sec:privacy}; TC-12 to TC-14 \\
Trust without blocking genuine emergencies & Identity verification; fail-open caller trust; audit trail & \cref{sec:caller-trust}; TC-06, TC-20, TC-21 \\
Hospital coordination with live vitals & Hospital selection, pre-notification, vitals stream, handoff & UC-43 to UC-46; TC-22, TC-23 \\
Arabic and English, both directions, phones and desktops & Localisation and right-to-left layout in both clients; bilingual notifications & FR-PLT-01, FR-NOT-02; TC-28 \\
\end{tabular}
\end{table}

\section{Limitations}

An honest assessment must also record what the system does not yet do.

\begin{itemize}
  \item \ProjectName{} does not connect to national emergency telephone numbers or to existing dispatch systems; it is
        designed to operate alongside them.
  \item Only an Android build of the mobile app exists, although the code base is cross-platform.
  \item The WebSocket channel is restricted by origin but individual subscriptions are not yet authorised by role.
  \item The rate limiter is held in memory in each instance, so running several backend instances would multiply the limits.
  \item Automated test coverage is concentrated on the rule classes; the end-to-end scenarios are run manually.
  \item The AI model service and its backend client are in place, but no clinical model has yet been trained and attached
        to the vitals stream.
\end{itemize}

\section{Future Work}\label{sec:future-work}

The improvements below follow directly from the limitations and from the extension points already built into the platform
(\cref{tab:future-work}).

{\tabfont\rowband
\begin{longtable}{@{}>{\raggedright\arraybackslash}p{\dimexpr(\linewidth-4\tabcolsep)*289/1000\relax}>{\raggedright\arraybackslash}p{\dimexpr(\linewidth-4\tabcolsep)*711/1000\relax}@{}}
\caption{Planned improvements}\label{tab:future-work}\\
\rowcolor{najdanavy}\textcolor{white}{\sffamily\bfseries\small Improvement} & \textcolor{white}{\sffamily\bfseries\small Description}\\
\endfirsthead
\rowcolor{najdanavy}\textcolor{white}{\sffamily\bfseries\small Improvement} & \textcolor{white}{\sffamily\bfseries\small Description}\\
\endhead
\bottomrule\endlastfoot
\textbf{Clinical vitals model} & Connect the AI model service (registry, contract and backend client already exist) to the transfer vitals stream so that hospitals see an acuity score beside each vitals record, with shadow-mode evaluation before it is shown.\\
\textbf{Wearable and sensor reports} & Accept automatically generated reports from consenting devices; the incident source field already reserves a sensor-device value.\\
\textbf{Channel-level WebSocket authorisation} & Authenticate STOMP connections and authorise each subscription by role and ownership, in addition to the current origin restriction.\\
\textbf{Distributed rate limiting} & Move the in-memory limiter to a shared store so limits hold across several API instances.\\
\textbf{Server-side unit ranking by road ETA} & Move ranking from the web console to the API using routed travel times, and expose it to other clients.\\
\textbf{iOS distribution} & Build and distribute the existing cross-platform app for iOS.\\
\textbf{Deeper automated testing} & Add service-level integration tests against a containerised PostgreSQL, and end-to-end browser tests for the console workflows.\\
\textbf{Operational analytics} & Response-time, coverage and false-report dashboards for administrators.\\
\textbf{Richer incident handling} & Sub-categories, multi-agency coordination and two-way voice or video with the caller.\\
\textbf{Offline maps and guidance packs} & Cache map tiles and guidance for areas with poor coverage.\\
\end{longtable}}


\section{Lessons Learned}

Several lessons stand out. First, \emph{decide the rules before the screens}: most of the difficult bugs of emergency
software are rule bugs --- races, retries, consent --- and writing the rules as small pure classes made them testable and
discussable. Second, \emph{an architecture proportionate to the team} matters more than fashionable architecture: a modular
monolith with one separate service gave strong boundaries without the cost of distribution. Third, \emph{degrade, do not
fail}: wrapping every external dependency so that it disables only its own feature kept the core workflow available.
Finally, \emph{documentation is part of the product}: keeping requirements, diagrams, data dictionary and traceability in step
with the code made the system far easier to reason about, to review and to hand over.

\section{Final Remarks}

\ProjectName{} demonstrates that a small team, using modern engineering practice and managed cloud services, can build a coherent
emergency-response platform that connects citizens, dispatchers, field crews and hospitals in a single, auditable workflow.
It is a foundation that can be extended with clinical models, additional channels and deeper integration with existing
emergency services.
